% ==========================================================================
% ch12.tex — dihasilkan otomatis dari website/chapters/ch12.qmd
% Boleh disunting manual. 'sync' hanya menimpa bila .qmd berubah DAN .tex
% belum disunting; kalau keduanya berubah, dilewati (pakai --force).
% qmd-sha1: f01f4d5e99a6
% tex-sha1: 3a4b9fd4437f
% ==========================================================================
% >>>>> ISI OTOMATIS (jangan hapus baris ini) >>>>>
\chapter{Citra \& Audio}\label{citra-audio}

\section{Citra: Normalisasi dan Augmentasi}\label{citra-normalisasi-dan-augmentasi}

Data citra digital pada dasarnya merupakan matriks numerik spasial berdimensi tinggi. Sebuah gambar berwarna tersusun atas kisi dua dimensi yang merekam tiga kanal warna (RGB), dengan setiap titik piksel menyimpan nilai intensitas dari 0 hingga 255. Sebelum arsitektur \emph{computer vision} dapat mengekstraksi representasi polanya, data mentah ini harus distandarkan melalui dua mekanisme utama: normalisasi dan augmentasi.

\textbf{Normalisasi Citra}
Normalisasi menstandarkan rentang nilai piksel untuk memastikan gradien algoritma \emph{machine learning} tetap stabil saat konvergensi. Praktik industri modern umumnya menerapkan standardisasi \emph{Z-score} per kanal warna melalui persamaan \( x_{c}^{\text{norm}} = \frac{x_c - \mu_c}{\sigma_c} \). Dalam rumusan tersebut, \(x_c\) mewakili intensitas piksel asli pada kanal warna spasial \(c\) (misalnya merah, hijau, atau biru), \(\mu_c\) menyatakan rata-rata intensitas piksel, dan \(\sigma_c\) merepresentasikan simpangan bakunya pada kanal tersebut. Untuk menjaga integritas \emph{pipeline}, parameter rata-rata dan simpangan baku wajib dihitung secara eksklusif dari himpunan data latih. Konstanta ini kemudian dikunci dan diaplikasikan pada data pengujian. Apabila menggunakan \emph{pretrained model}, praktisi harus menormalisasi citra mengikuti distribusi statistik asal model tersebut, seperti konstanta rata-rata standar ImageNet \passthrough{\lstinline![0.485, 0.456, 0.406]!}.

\textbf{Augmentasi Data}
Augmentasi menyisipkan transformasi sintetis di dalam \emph{pipeline} untuk mencegah model menghafal tata letak piksel (\emph{overfitting}). Transformasi ini secara khusus hanya dijalankan pada fase pelatihan. Teknik augmentasi mencakup manipulasi geometris dan visual dasar, seperti pemotongan acak, pembalikan horizontal, distorsi prespektif, dan injeksi \emph{noise} warna. Selain itu, kerangka kerja modern menggunakan kebijakan otomatis (\emph{AutoAugment} atau \emph{RandAugment}) yang secara acak mengundi kombinasi transformasi dari kumpulan prosedur operasi standar. Terdapat juga pendekatan tingkat lanjut seperti augmentasi lintas sampel (\emph{Mixup} atau \emph{CutMix}) yang menggabungkan potongan citra dari kelas berbeda untuk melatih model memprediksi rasio probabilitas gabungan, sehingga model difokuskan pada morfologi lokal dibandingkan latar belakang global.

Kedua tahapan prapemrosesan ini menentukan apakah representasi fitur dapat mengekstraksi objek dengan konsisten atau terdistorsi oleh variasi pencahayaan.

\section{Fitur Rancangan Tangan (HOG) dan Keterbatasannya}\label{fitur-rancangan-tangan-hog-dan-keterbatasannya}

Sebelum representasi mendalam digunakan secara luas, pemahaman citra bertumpu pada representasi rancangan manusia. Pendekatan ini menghindari penggunaan piksel mentah yang sensitif terhadap perubahan cahaya atau pergeseran posisi. Pergeseran pencahayaan dapat mengubah susunan angka piksel secara drastis, meskipun objek yang ditangkap tetap sama. Ekstraksi fitur difokuskan pada perhitungan lokal yang memanfaatkan operasi matematika untuk mendeteksi tepian atau sudut pada area kecil gambar. Hasil perhitungan tersebut mengubah matriks piksel menjadi vektor fitur yang konsisten terhadap gangguan visual lingkungan.

Contoh paling menonjol dari pendekatan ini adalah algoritma \emph{Histogram of Oriented Gradients} (HOG). Algoritma ini membagi citra ke dalam kisi-kisi blok berukuran kecil dan menghitung arah perubahan intensitas warna (\emph{gradient}). Dalam kalkulasi ini, gradien citra dihitung berdasarkan perbedaan intensitas antara piksel yang bersebelahan. Jika \(I(x,y)\) mewakili intensitas piksel pada koordinat \((x,y)\), komponen gradien horizontal (\(G_x\)) dan vertikal (\(G_y\)) dirumuskan sebagai \(G_x = I(x+1, y) - I(x-1, y)\) dan \(G_y = I(x, y+1) - I(x, y-1)\). Variabel \(G_x\) mengukur perubahan intensitas secara horizontal untuk memetakan tepian vertikal, sedangkan \(G_y\) mengukur perubahan vertikal untuk memetakan tepian horizontal. Dari kedua nilai tersebut, algoritma HOG menghitung magnitudo dan orientasi arah gradien untuk menyusun siluet kontur fisik objek.

\begin{figure}[htbp]\centering
\includegraphics[width=\linewidth,height=0.72\textheight,keepaspectratio]{ch12-fig-1}
\caption{Ekstraksi Fitur Histogram of Oriented Gradients (HOG) pada Citra}
\label{fig:ch12-fig-1}
\end{figure}

Pada kasus deteksi pejalan kaki, model tidak mengevaluasi warna pakaian target. Ekstraktor HOG memetakan siluet target dengan mengidentifikasi garis vertikal tubuh serta lengkungan kepala melalui nilai gradien. Perhitungan tersebut kemudian dikelompokkan ke dalam blok-blok spasial berukuran tetap. Algoritma merangkum dominasi arah gradien tersebut ke dalam sebuah histogram distribusi orientasi lokal. Representasi spasial ini bertindak sebagai fitur akhir, sehingga target terdeteksi secara konsisten di dalam ruang gradien HOG tanpa terpengaruh variasi visual pakaian.

Pendekatan ekstraksi manual ini memiliki keterbatasan struktural. Pembuatan algoritma menuntut pemahaman domain spesifik mengenai manipulasi citra dan proses rekayasa parametrik. Parameter sistem seperti dimensi sel, ukuran blok, dan ambang batas intensitas harus dikonfigurasi secara manual untuk target deteksi yang berbeda. Selain itu, fitur rancangan tangan tidak memetakan konteks tingkat tinggi. HOG memetakan batas siluet objek dengan baik, tetapi tidak mampu menyerap konteks semantik, seperti membedakan objek fisik dari gambar dua dimensi. Keterbatasan ini memfasilitasi transisi penggunaan representasi yang dipelajari secara langsung dari data pelatihan.

\section{\texorpdfstring{CNN \emph{Feature Extractor} dan \emph{Pretrained Model}}{CNN Feature Extractor dan Pretrained Model}}\label{cnn-feature-extractor-dan-pretrained-model}

Pemrosesan citra berkembang secara signifikan melalui arsitektur \emph{Convolutional Neural Network} (CNN) yang memfasilitasi ekstraksi representasi melalui pembelajaran mesin. Algoritma ini memetakan pola visual dari data pelatihan tanpa menggunakan aturan matematis yang statis. CNN bekerja dengan menggeser sekumpulan filter melintasi area citra untuk mengekstraksi berbagai fitur. Operasi konvolusi dasar dirumuskan sebagai \( \mathbf{Y}_{i,j} = \sum_{m} \sum_{n} \mathbf{W}_{m,n} \cdot \mathbf{X}_{i+m, j+n} \). Dalam perumusan tersebut, \(\mathbf{Y}\) melambangkan \emph{feature map} keluaran, \(\mathbf{W}\) merupakan matriks filter pembawa bobot yang dioptimalkan selama pelatihan, dan \(\mathbf{X}\) merepresentasikan area piksel citra masukan.

Melalui operasi konvolusi berulang, CNN membangun pemahaman visual secara hierarkis. Pada lapisan awal, filter merespons komponen spasial sederhana seperti garis vertikal, sudut, lengkungan, dan gradasi warna. Pada lapisan menengah, jaringan menyatukan komponen tersebut menjadi struktur tekstur atau fitur spesifik objek. Pada lapisan dalam, fitur parsial dikonsolidasi menjadi representasi semantik tingkat tinggi yang membedakan objek secara utuh. Kemampuan pembelajaran berjenjang ini mendukung implementasi \emph{transfer learning}. CNN yang telah dilatih dengan jutaan citra memiliki representasi spasial yang universal. Model dasar ini bertindak sebagai \emph{pretrained model}, di mana kemampuan ekstraksinya digunakan kembali untuk mengevaluasi domain data baru.

\begin{figure}[htbp]\centering
\includegraphics[width=\linewidth,height=0.72\textheight,keepaspectratio]{ch12-fig-2}
\caption{Arsitektur CNN Menghilangkan Lapisan Klasifikasi Akhir Menjadi Feature Extractor}
\label{fig:ch12-fig-2}
\end{figure}

Dalam implementasinya, varian arsitektur seperti ResNet dapat dimodifikasi dengan menghilangkan lapisan klasifikasi akhir. Sisa jaringan dialihfungsikan sebagai \emph{feature extractor} murni. Ketika citra medis atau foto industri dimasukkan ke dalam jaringan ini, model tidak memprediksi probabilitas kelas, melainkan memproduksi vektor fitur yang memuat representasi visual citra. Vektor inilah yang kemudian dialirkan ke model prediktif sederhana untuk mendeteksi anomali. Pendekatan ini menghasilkan representasi dengan akurasi tinggi tanpa mengharuskan sistem melatih parameter dari awal.

Meskipun model berbasis \emph{Vision Transformer} (ViT) menjadi standar alternatif yang populer, CNN tetap dipertahankan pada lini komputasi industri. Evaluasi empiris menunjukkan bahwa arsitektur CNN seperti ConvNeXt v2 memiliki kinerja yang sebanding atau mengungguli ViT pada himpunan data berukuran menengah. Struktur operasi konvolusi memungkinkan model mengekstraksi informasi visual dengan efisien saat ketersediaan data pelatihan sangat terbatas.

\section{Image Embedding dan Patch}\label{image-embedding-dan-patch}

Arsitektur \emph{Vision Transformer} (ViT) menyediakan pendekatan alternatif yang memproses seluruh area citra sebagai urutan sekuensial. Langkah awal metode ini membagi citra menjadi kumpulan \emph{patch} persegi yang diletakkan berdampingan tanpa tumpang tindih. Setiap \emph{patch} kemudian diratakan dan diproses sebagai sebuah \emph{token} independen, yang memiliki peran identik dengan representasi kata dalam model bahasa. Pemotongan area spasial ini dirumuskan melalui persamaan \( N = \frac{H \times W}{P^2} \). Dalam persamaan tersebut, \(N\) menyatakan total \emph{patch} keluaran, \(H\) dan \(W\) mewakili dimensi tinggi dan lebar citra asli, serta \(P\) merepresentasikan ukuran sisi setiap \emph{patch}. Sebagai contoh, citra beresolusi 224×224 piksel dengan ukuran \emph{patch} 16×16 piksel akan direpresentasikan menjadi 256 \emph{patch} terpisah.

Komponen komputasi \emph{attention} mengevaluasi keterkaitan seluruh \emph{token} tersebut secara bersamaan. Mekanisme ini memfasilitasi identifikasi korelasi antara objek yang terpisah secara spasial dalam satu gambar.

\begin{figure}[htbp]\centering
\includegraphics[width=\linewidth,height=0.72\textheight,keepaspectratio]{ch12-fig-3}
\caption{Pemotongan Citra Menjadi Kisi Patch dan Mekanisme Attention pada Vision Transformer}
\label{fig:ch12-fig-3}
\end{figure}

Pemrosesan relasi antarkepingan ini mengonsolidasi makna semantik visual ke dalam ruang kontinu tunggal yang disebut \emph{image embedding}. Representasi ini memetakan konsep objek secara utuh. Jika dua citra menampilkan kelas objek yang sama, vektor \emph{embedding} dari keduanya akan dipetakan berdekatan secara matematis tanpa dipengaruhi oleh variasi latar belakang visual.

Dalam praktik rekayasa fitur modern, terdapat beberapa keluarga \emph{pretrained model} yang menjadi standar ekstraksi citra. Model keluarga CLIP dan SigLIP dilatih menggunakan jutaan pasangan citra dan teks, dengan bagian \emph{encoder} visual dari SigLIP sering digunakan sebagai pengekstraksi fitur serbaguna karena memiliki rentang representasi yang lebih umum dibanding arsitektur konvensional. Alternatif lainnya adalah DINOv2, arsitektur \emph{self-supervised learning} yang dioptimalkan untuk mengekstraksi struktur spasial. DINOv2 memproduksi \emph{token} CLS yang merangkum \emph{image embedding} global, bersama fitur tingkat \emph{patch} yang mendukung resolusi tinggi untuk memfasilitasi tugas prediksi padat seperti segmentasi spasial. Untuk aplikasi tingkat lanjut, fitur direpresentasikan menggunakan pendekatan multi-skala dan hibrida. Arsitektur modern seperti ConvNeXt v2 menggabungkan hierarki CNN dengan protokol pelatihan transformer untuk mengompilasi fitur yang merekam detail lokal dan konteks global sekaligus.

Pendekatan sekuensial ini menyelaraskan pemrosesan data lintas modalitas, di mana representasi visual dan tekstual disatukan menggunakan struktur komputasi yang seragam pada arsitektur pembelajaran multimodal.

\section{Audio: Waveform, Spektrogram, MFCC, dan Chroma}\label{audio-waveform-spektrogram-mfcc-dan-chroma}

Data audio pada dasarnya merupakan sinyal \emph{waveform} mentah berdimensi tunggal yang mencatat perubahan amplitudo tekanan udara seiring waktu. Pada setiap titik temporal, berbagai sumber suara berakumulasi menjadi satu nilai amplitudo. Mengingat sinyal ini menuntut pemisahan komponen akustik, pendekatan klasik merancang transformasi untuk mengonversi sinyal dari domain waktu ke domain frekuensi melalui spektrogram. Transformasi ini memetakan fluktuasi satu dimensi menjadi representasi visual dua dimensi. Rentang waktu diposisikan pada sumbu horizontal, frekuensi pada sumbu vertikal, dan amplitudo direpresentasikan melalui intensitas piksel. Perubahan format menjadi citra spektrum memfasilitasi penerapan algoritma \emph{computer vision} pada data audio.

Untuk memodelkan persepsi pendengaran biologis, dirancang representasi frekuensi yang mengikuti skala Mel. Pendengaran manusia beroperasi secara non-linear, menunjukkan sensitivitas tinggi terhadap pergeseran nada pada frekuensi rendah dan penurunan sensitivitas pada frekuensi tinggi. Korelasi antara frekuensi asli dengan skala Mel dinyatakan melalui persamaan \( m = 2595 \log_{10} \left( 1 + \frac{f}{700} \right) \). Dalam perumusan tersebut, \(m\) menyatakan frekuensi pada skala Mel dan \(f\) menyatakan frekuensi aktual dalam satuan Hertz.

Berdasarkan konversi persepsi tersebut, terdapat beberapa teknik ekstraksi fitur spesifik untuk pemrosesan audio. Mel-spektrogram memetakan sumbu frekuensi ke skala Mel sambil mempertahankan dimensi spasial komponen suara, menjadikannya format standar bagi pengklasifikasi audio berbasis CNN. Teknik \emph{Mel-Frequency Cepstral Coefficients} (MFCC) menghasilkan representasi yang mengekstraksi selubung spektrum suara dari Mel-spektrogram. MFCC umumnya diimplementasikan ketika data latih berjumlah kecil untuk pemodelan klasik, atau pada arsitektur \emph{edge deployment} dengan kendala memori. Untuk analisis musik, \emph{chroma feature} meringkas energi frekuensi ke dalam dua belas kelas nada dasar. Teknik ini mengeliminasi informasi oktaf untuk menonjolkan identitas harmoni, sehingga mendukung inferensi progresi kord dan klasifikasi genre lagu. 

Teknik-teknik ekstraksi ini merepresentasikan fitur berbasis rancangan manusia yang digunakan untuk mereduksi dimensi sekaligus mengisolasi pola sinyal akustik sebelum dievaluasi oleh sistem prediktif.

\section{Pretrained Audio Encoder: Model Raw-Waveform}\label{pretrained-audio-encoder-model-raw-waveform}

Pemrosesan audio mutakhir menggunakan representasi yang dipelajari mesin untuk mengolah \emph{raw waveform} secara langsung. Arsitektur \emph{state-of-the-art} mengekstrak pola akustik dari fluktuasi amplitudo tanpa menggunakan transformasi domain frekuensi statis. Model \emph{raw-waveform} seperti wav2vec 2.0 dan HuBERT memproses deret amplitudo 1D menggunakan tumpukan konvolusi 1D. Alih-alih mengekstraksi fitur diskrit melalui konversi tetap, model ini menggunakan matriks bobot yang dioptimalkan dari data pelatihan untuk memetakan gelombang suara menjadi matriks fitur kontinu.

Peningkatan metrik arsitektur ini didukung oleh metode \emph{self-supervised learning} (SSL). Selama pelatihan, model mengevaluasi puluhan ribu jam rekaman wicara tanpa label. Sistem menerapkan penyembunyian (\emph{masking}) pada sebagian segmen deret representasi, dan model dioptimalkan untuk memprediksi fitur dari porsi yang disembunyikan. Tugas rekonstruksi ini memungkinkan model mengenali hierarki fonetik secara otonom.

Implementasi rekayasa fitur saat ini bergeser pada penggunaan \emph{pretrained model} untuk menghasilkan vektor representatif. Keluarga model wav2vec 2.0 dan HuBERT bertindak sebagai standar untuk tugas pengenalan ucapan (\emph{speech recognition}), dengan iterasi besar seperti Wav2Vec2-BERT 2.0 yang dilatih menggunakan jutaan jam audio. Model WavLM mengembangkan arsitektur HuBERT menggunakan mekanisme pencampuran ucapan (\emph{utterance mixing}) untuk menahan retensi identitas pembicara, yang diutamakan pada tugas \emph{speaker verification}. Pada klasifikasi alternatif, status tersembunyi (\emph{hidden state}) dari blok akhir \emph{encoder} model Whisper digunakan sebagai pengekstraksi fitur semantik. Selain arsitektur wicara, terdapat model spesifik-domain seperti AudioMAE dan BEATs yang dikalibrasi untuk pengenalan kejadian suara lingkungan (\emph{sound event detection}), serta model MERT untuk mengekstraksi fitur kemusikan.

Pada tahapan pra-pemrosesan, rentang sampel \emph{waveform} distandarkan durasinya melalui pemotongan atau penambahan bantalan (\emph{padding}), dan laju sampelnya (\emph{sampling rate}) diseragamkan ke batas spesifikasi arsitektur sebelum dievaluasi oleh \emph{encoder}. Keluaran dari tahap ini kemudian difungsikan sebagai masukan bagi klasifikasi hilir.

\section{\texorpdfstring{\emph{Audio Embedding} dan Agregasi Sepanjang Waktu}{Audio Embedding dan Agregasi Sepanjang Waktu}}\label{audio-embedding-dan-agregasi-sepanjang-waktu}

Rekaman audio memuat dimensi temporal melintasi deret waktu. Pemrosesan rekaman melalui ekstraktor fitur menghasilkan rentetan matriks observasi dalam orde milidetik. Karakteristik sekuensial ini menuntut mekanisme agregasi, mengingat pengklasifikasi hilir (\emph{downstream classifier})—seperti model deteksi kelas suara—membutuhkan masukan dengan dimensi statis yang tidak bergantung pada panjang durasi asli. Deretan fitur tingkat \emph{frame} dikompresi menjadi vektor representatif tunggal yang disebut \emph{audio embedding}.

Pemadatan (\emph{pooling}) representasi temporal ini dieksekusi melalui beberapa mekanisme agregasi. Metode \emph{average pooling} menghitung nilai rata-rata dari seluruh kepingan \emph{frame} dan diterapkan untuk representasi sinyal kontinu dengan profil statis, seperti dengung mesin. Metode \emph{max pooling} menyeleksi nilai aktivasi matriks tertinggi untuk mempertahankan sinyal akustik impulsif yang berdurasi singkat, sehingga mencegah pengenceran nilai oleh \emph{frame} hening. Pada arsitektur yang lebih kompleks, mekanisme \emph{attention} dinamis memindai rekaman secara adaptif dan memberikan bobot fitur yang lebih besar pada kepingan waktu yang memuat relevansi semantik.

Integrasi ini menghasilkan \emph{embedding} seragam yang mengonversi variabilitas temporal menjadi representasi geometri spasial yang stabil. Representasi berdimensi tetap ini mendukung kompatibilitas format masukan untuk pengklasifikasi standar dan memfasilitasi fusi fitur lintas modalitas pada arsitektur pembelajaran multimodal.

\section{Studi Kasus: Fitur Rancangan Tangan vs.~Representasi Mendalam}\label{studi-kasus-fitur-rancangan-tangan-vs.-representasi-mendalam}

Notebook praktikum pada bab ini mendemonstrasikan perbandingan antara representasi yang dirancang manusia dengan representasi yang dipelajari mesin dalam menangani variasi observasi.

Pada domain visual untuk klasifikasi kendaraan, pengklasifikasi \emph{baseline} menggunakan representasi HOG yang memetakan distribusi kemiringan tepi spasial. Evaluasi menunjukkan bahwa oklusi parsial pada objek merusak siluet fitur HOG, sehingga akurasi identifikasi menurun drastis. Sebaliknya, \emph{pretrained model} berbasis CNN mampu mengidentifikasi komponen semantik dari fitur lokal seperti lampu atau roda. Arsitektur ini dapat mempertahankan kinerja klasifikasi pada kondisi oklusi tanpa memerlukan visibilitas siluet objek yang utuh.

Pada klasifikasi domain akustik, ekstraktor klasik MFCC memampatkan rentang spektrum ke dalam representasi berdimensi rendah. Proses pemadatan ini menyebabkan sinyal derau lingkungan melebur dengan frekuensi suara target. Sebaliknya, \emph{audio embedding} yang diekstraksi dari model \emph{raw-waveform} mengisolasi korelasi ruang frekuensi dan merepresentasikan pola suara target dari spektrum latarnya.

{[}GAMBAR 12.4: Heatmap - Visualisasi MFCC yang memampatkan rentang frekuensi suara, menunjukkan peleburan derau lingkungan dengan sinyal target{]}

Penurunan resolusi pada MFCC terkait erat dengan mekanisme transformasinya yang menggunakan persamaan skala Mel \( m = 2595 \log_{10} \left(1 + \frac{f}{700}\right) \). Fungsi logaritmik dalam persamaan ini secara inheren mereduksi resolusi pada pita frekuensi tinggi. Sebagai konsekuensinya, pada observasi yang mengandung distorsi derau, variasi sinyal suara target tidak dapat dipisahkan secara optimal oleh arsitektur pengklasifikasi hilir.

Implementasi rekayasa fitur pada sekuens berkelanjutan memiliki kerentanan terhadap kebocoran data (\emph{data leakage}). Skema partisi (\emph{split}) data tingkat klip—di mana satu sesi rekaman utuh dipotong dan didistribusikan secara acak ke set pelatihan dan pengujian—akan memvalidasi model pada pola akustik lingkungan yang persis sama. Partisi ini memicu model untuk sekadar mempelajari korelasi latar belakang. Praktik rekayasa standar mensyaratkan pembagian pada tingkat sumber rekaman (\emph{source-level split}). Pada skema ini, seluruh klip dari sesi perekaman fisik tertentu dialokasikan secara utuh ke dalam himpunan pelatihan atau himpunan pengujian tanpa perpotongan.
